\section{Architecture}
\label{sec:architecture}

\alaya{} is organized around three memory stores connected by a Hebbian graph
overlay, with five lifecycle operations that manage memory health over time.
\Cref{fig:architecture} illustrates the overall structure.

\begin{figure}[t]
\centering
\begin{tikzpicture}[
    store/.style={draw, rounded corners, minimum width=2.2cm, minimum height=1cm, font=\small},
    op/.style={draw, dashed, rounded corners, minimum width=1.8cm, minimum height=0.7cm, font=\scriptsize, fill=gray!10},
    arr/.style={-{Stealth[length=2mm]}, thick},
]
    % Stores
    \node[store, fill=blue!15] (ep) at (0, 0) {Episodic};
    \node[store, fill=green!15] (sem) at (3.2, 0) {Semantic};
    \node[store, fill=orange!15] (imp) at (6.4, 0) {Implicit};

    % Graph overlay
    \node[draw, rounded corners, minimum width=7.5cm, minimum height=0.7cm,
          fill=yellow!10, font=\small] (graph) at (3.2, -1.5) {Hebbian Graph Overlay};

    % Lifecycle ops
    \node[op] (cons) at (1.6, 1.5) {Consolidate};
    \node[op] (trans) at (3.2, 2.3) {Transform};
    \node[op] (forg) at (4.8, 1.5) {Forget};
    \node[op] (perf) at (0, 2.3) {Perfume};
    \node[op] (onto) at (6.4, 2.3) {Ontology};

    % Arrows
    \draw[arr] (ep) -- (cons);
    \draw[arr] (cons) -- (sem);
    \draw[arr] (trans) -- (sem);
    \draw[arr] (forg) -- (sem);
    \draw[arr] (perf) -- (imp);
    \draw[arr] (onto) -- (sem);

    % Graph connections
    \draw[arr, gray] (ep) -- (graph);
    \draw[arr, gray] (sem) -- (graph);
    \draw[arr, gray] (imp) -- (graph);
\end{tikzpicture}
\caption{\alaya{} architecture. Three stores hold different memory types;
the Hebbian graph overlay connects them with weighted, use-adaptive links.
Five lifecycle operations (dashed) manage memory beyond storage and
retrieval.}
\label{fig:architecture}
\end{figure}

\subsection{Memory Stores}
\label{sec:stores}

\paragraph{Episodic store (hippocampus)}
Records individual conversational turns with full context: content, speaker
role, session identifier, timestamp, topic tags, entities, sentiment, and
optional embedding vector. Episodes are indexed for full-text search (FTS5)
and vector similarity. Following the CLS model, the episodic store is
high-fidelity and append-only; management occurs through consolidation and
forgetting, not in-place mutation.

\paragraph{Semantic store (neocortex)}
Holds stable knowledge extracted during consolidation: facts
(\emph{``User programs in Rust''}), relationships (\emph{``Alice manages
Bob''}), events (\emph{``User presented at RustConf 2025''}), and concepts
(\emph{``Functional programming emphasizes immutability''}). Each semantic
node carries a confidence score, corroboration count (how many independent
episodes support it), source episode links, and optional category assignment.
Semantic nodes are the primary substrate for long-term memory.

\paragraph{Implicit store (v\=asan\=a)}
Models preferences that emerge from accumulated interaction rather than
explicit extraction. Raw \emph{impressions}---observations about user
behavior (\emph{``User prefers concise responses''})---accumulate over many
sessions. When impressions converge (five or more similar observations), a
\emph{preference} crystallizes with a confidence score. This models the
Yogac\=ara concept of v\=asan\=a (perfuming): preferences are not extracted
from any single interaction but emerge gradually from the ``incense'' of
repeated experience.

\subsection{Hebbian Graph Overlay}
\label{sec:graph}

All three stores are connected by a shared graph with six link types:

\begin{itemize}[nosep]
    \item \textbf{Temporal:} Sequential relationships between episodes in a
    session.
    \item \textbf{Topical:} Shared topic tags between any two nodes.
    \item \textbf{Entity:} Shared named entities.
    \item \textbf{Causal:} Explicit causal relationships identified during
    consolidation.
    \item \textbf{CoRetrieval:} Created when two nodes are returned in the
    same query result, modeling the Hebbian principle that co-activated
    memories strengthen their association.
    \item \textbf{MemberOf:} Category membership links created by the
    emergent ontology system (\Cref{sec:ontology}).
\end{itemize}

Link weights evolve through use: co-retrieval strengthens weights (LTP),
while the passage of time without co-activation allows weights to decay
(LTD). This produces emergent small-world topology in the graph, where
frequently co-accessed memories form tightly connected clusters while rarely
used associations weaken. The graph is used during retrieval for spreading
activation---a query activates seed nodes, and activation propagates through
weighted links with distance-based decay, surfacing indirectly related
memories that pure similarity search would miss.

\subsection{Lifecycle Operations}
\label{sec:lifecycle}

\paragraph{Consolidation (\texttt{consolidate()})}
Implements CLS-style interleaved replay. Recent unconsolidated episodes are
batched and submitted to a \texttt{ConsolidationProvider} (typically an LLM)
that extracts structured semantic nodes---facts, relationships, events, and
concepts. The provider returns typed nodes with confidence scores and source
episode references. \alaya{} stores these in the semantic store, creates
Hebbian links to source episodes, initializes Bjork dual-strength values,
and attempts category assignment (\Cref{sec:ontology}). The provider
interface is pluggable: applications supply their own LLM integration,
keeping \alaya{} independent of any specific model or API.

In Yogac\=ara terms, consolidation implements \emph{vip\=aka}
(heterogeneous ripening): the output is qualitatively different from the
input. Episodes are specific, contextual, and time-bound; semantic nodes
are decontextualized, typed, and reusable.

\paragraph{Transformation (\texttt{transform()})}
Periodic reorganization of the semantic store, implementing
\emph{\=a\'sraya-par\=av\d{r}tti} (transformation of the basis). Three
sub-operations run in sequence:
\begin{enumerate}[nosep]
    \item \textbf{Deduplication:} Semantic nodes with high embedding
    similarity and matching type are merged, combining corroboration counts
    and source links.
    \item \textbf{Contradiction resolution:} Nodes flagged as contradictory
    (opposing confidence signals from corroboration) are resolved by
    retaining the more strongly supported version.
    \item \textbf{Category discovery:} Uncategorized semantic nodes are
    clustered to discover emergent categories (\Cref{sec:ontology}).
\end{enumerate}

\paragraph{Forgetting (\texttt{forget()})}
Implements Bjork dual-strength decay. Retrieval strength decays as a
function of time since last access; storage strength is monotonically
non-decreasing. When retrieval strength drops below a configurable
threshold, the memory becomes effectively inaccessible---still stored but
unlikely to be retrieved. The \texttt{forget()} operation can purge
memories below this threshold, modeling the Yogac\=ara principle that
unactivated seeds weaken over time. During query, retrieval-induced
forgetting suppresses competitors: memories sharing retrieval cues with
the accessed memory have their retrieval strength reduced, implementing
the competitive dynamics that Anderson~\etal~\cite{anderson1994}
demonstrated in human memory.

\paragraph{Perfuming (\texttt{perfume()})}
Implements v\=asan\=a (gradual impression accumulation). Each interaction
can generate impressions---observations about user behavior, preferences,
or communication style. Impressions accumulate in the implicit store. When
five or more similar impressions converge, a preference crystallizes. This
gradual process avoids the brittleness of one-shot preference extraction:
a user mentioning Python once does not generate a ``prefers Python''
preference, but mentioning it across fifty sessions does.

\paragraph{Emergent ontology}
The newest lifecycle operation, detailed in \Cref{sec:ontology}. Categories
emerge from clustering patterns in the semantic store, providing a
lightweight taxonomic structure without manual curation or supervision.

\subsection{Retrieval Pipeline}
\label{sec:retrieval}

\alaya{} implements Modular RAG~\cite{gao2023rag} with four retrieval
signals fused via Reciprocal Rank Fusion (RRF)~\cite{cormack2009rrf}:

\begin{enumerate}[nosep]
    \item \textbf{BM25 (sparse):} Full-text search via SQLite FTS5,
    effective for exact term matches and proper nouns.
    \item \textbf{Vector (dense):} Cosine similarity over embedding
    vectors, effective for semantic similarity.
    \item \textbf{Graph (associative):} Spreading activation from seed
    nodes through the Hebbian graph overlay, effective for indirect
    associations.
    \item \textbf{Reranking:} Hebbian co-retrieval frequency and Bjork
    retrieval strength boost results that have been useful in similar past
    contexts.
\end{enumerate}

\noindent RRF combines ranked lists from each signal without requiring
score normalization: $\text{RRF}(d) = \sum_{r \in R} \frac{1}{k + r(d)}$,
where $r(d)$ is the rank of document $d$ in ranked list $r$ and $k$ is a
constant (default 60). This is robust to heterogeneous score distributions
across signals, which is critical when combining sparse, dense, and
graph-based scores.
